%! Author = lili
%! Date = 7/28/26
\section{Abgabe 1.1}

\subsection{HA 1.1.1 (Lösung vorhanden, nicht eingetragen)}

\subsection{HA 1.1.2 (Lösung vorhanden, nicht eingetragen)}

\subsection{HA 1.1.3 (Lösung vorhanden, nicht eingetragen)}
Wir drehen ein Glücksrad dreimal hintereinander.
Das Glücksrad ist in vier gleichgroße Viertel unterteilt, die mit 0, 1, 2 und 3 beschriftet sind.

\begin{auf}[Geben Sie einen geeigneten Ereignisraum $\Omega$ für dieses Experiment an, der es erlaubt, die folgenden Fragen zu beantworten.]
    LÖSUNG DA, NICHT EINGETRAGEN % TODO NOCH KEINE LÖSUNG
\end{auf}

\begin{auf}[Geben Sie eine für dieses Zufallsexperiment geeignete Wahrscheinlichkeitsfunktion $\mu$ an. Begründen Sie Ihre Wahl.]
    LÖSUNG DA, NICHT EINGETRAGEN % TODO NOCH KEINE LÖSUNG
\end{auf}

\begin{auf}[Definieren Sie Ereignisse $A$, $B$ und $C$ als Teilmengen von $\Omega$ derart, dass...]
    \begin{itemize}
        \item $A$ das Ereignis ist, dass die zweite Drehung eine 3 ergibt;
        \item $B$ das Ereignis ist, dass genau zwei der drei Drehungen die gleiche Zahl ergeben und das Maximum aller Ergebnisse 1 ist;
        \item $C$ das Ereignis ist, dass die Zahlen gemäß der Reihenfolge des Drehens eine strikt wachsende Zahlenfolge bilden.\footnote{D.\,h., das Ergebnis des zweiten Drehens ist echt größer als das Ergebnis des ersten Drehens, und das Ergebnis des dritten Drehens ist wiederum echt größer als das Ergebnis des zweiten Drehens.}
    \end{itemize}
    LÖSUNG DA, NICHT EINGETRAGEN % TODO NOCH KEINE LÖSUNG
\end{auf}

\begin{auf}[Berechnen Sie die Wahrscheinlichkeiten von $A$, $B$ und $C$ mit Hilfe der Wahrscheinlichkeitsfunktion $\mu$.]
    LÖSUNG DA, NICHT EINGETRAGEN % TODO NOCH KEINE LÖSUNG
\end{auf}